\documentclass[9pt,landscape,a4paper]{article}
\usepackage[utf8]{inputenc}
\usepackage[T1]{fontenc}
\usepackage{geometry}
\usepackage{multicol}
\usepackage{amsmath,amssymb}
\usepackage{enumitem}
\usepackage{graphicx}
\usepackage{xcolor}
\usepackage{listings}
\usepackage{booktabs}
\usepackage{titlesec}
\usepackage{fancyhdr}
\usepackage{hyperref}
\usepackage{CJKutf8}

% Page layout
\geometry{top=0.7cm, bottom=0.7cm, left=0.7cm, right=0.7cm}
\pagestyle{empty}
\setlength{\parindent}{0pt}
\setlength{\parskip}{1.5pt}
\setlength{\columnsep}{7pt}

% Section formatting
\titleformat{\section}{\normalsize\bfseries\color{blue!70!black}}{}{0em}{}[\titlerule]
\titleformat{\subsection}{\small\bfseries\color{blue!50!black}}{}{0em}{}
\titlespacing{\section}{0pt}{4pt}{2pt}
\titlespacing{\subsection}{0pt}{3pt}{1pt}

% Code listing style
\lstset{
    language=Python,
    basicstyle=\ttfamily\fontsize{6.5}{7.5}\selectfont,
    keywordstyle=\color{blue!70!black}\bfseries,
    stringstyle=\color{red!60!black},
    commentstyle=\color{green!50!black}\itshape,
    numberstyle=\tiny\color{gray},
    backgroundcolor=\color{gray!8},
    frame=single,
    framerule=0.3pt,
    rulecolor=\color{gray!50},
    breaklines=true,
    breakatwhitespace=false,
    tabsize=2,
    showstringspaces=false,
    aboveskip=2pt,
    belowskip=2pt,
    xleftmargin=2pt,
    xrightmargin=2pt,
}

\newcommand{\keybox}[1]{%
    \begin{center}
    \fcolorbox{blue!50!black}{blue!5}{%
        \parbox{0.95\linewidth}{\small #1}%
    }
    \end{center}
}

\newcommand{\zhint}[1]{{\scriptsize\color{gray!70!black}#1}}

\begin{document}
\begin{CJK}{UTF8}{gbsn}

\begin{center}
{\large\bfseries\color{blue!70!black} Chapter 13 -- Course Summary}\\
{\footnotesize IERG4320/IEMS5726 Data Science in Practice | Cheat Sheet}
\end{center}

\begin{multicols}{3}

% ==================== SECTION 1 ====================
\section{Exam Information}
\zhint{考试信息}

\textbf{Format} (3 types):
\begin{enumerate}[leftmargin=10pt,itemsep=0pt,topsep=0pt]
\item \textbf{Python Programming} (30\%): 3 problems -- data preprocessing, ML, visualization. Fill in the blanks.
\item \textbf{Short Questions} (20\%): many small problems not covered by scenario-based.
\item \textbf{Scenario-based Written} (50\%): 4 problems -- beginner to moderate interview questions.
\end{enumerate}

\textbf{Tips}: revise assignment Q1, lecture examples, UReply questions, prepare DS interview questions, understand course project design.

% ==================== SECTION 2 ====================
\section{Data Science Pipeline}
\zhint{数据科学流程：定义问题→收集→预处理→建模→解释→进一步分析}

\fbox{\parbox{0.95\linewidth}{\scriptsize Problem Def $\to$ Collection $\to$ Preprocessing $\to$ Model $\to$ Interpret $\to$ Further Analysis}}

\subsection{Data Preprocessing (8 Steps)}
\zhint{数据预处理8步}
\begin{enumerate}[leftmargin=10pt,itemsep=0pt,topsep=0pt]
\item Load libraries and dataset
\item Data wrangling \zhint{数据整理}
\item Check missing values \zhint{检查缺失值}
\item Handle categorical values \zhint{处理分类变量}
\item Handle NL, images, audio \zhint{处理文本/图像/音频}
\item Dimensionality reduction \zhint{降维}
\item Train/validation/test split \zhint{划分数据集}
\item Feature scaling \zhint{特征缩放}
\end{enumerate}

\subsection{4 Types of Data Analytics}
\zhint{四种数据分析}
\begin{itemize}[leftmargin=8pt,itemsep=0pt,topsep=0pt]
\item \textbf{Descriptive}: what happened? \zhint{描述性：发生了什么}
\item \textbf{Diagnostic}: why did it happen? \zhint{诊断性：为什么}
\item \textbf{Predictive}: what will happen? \zhint{预测性：会发生什么}
\item \textbf{Prescriptive}: what should we do? \zhint{规范性：应该怎么做}
\end{itemize}

\subsection{Basic ML Recap}
\zhint{机器学习基础回顾}

\textbf{4 data formats}: numerical, categorical, ordinal, text/image/audio.\\
\textbf{4 ML problem types}: classification, regression, clustering, reinforcement learning.\\
\textbf{3 algorithm types}: supervised, unsupervised, reinforcement.\\
\textbf{3 dataset types}: training, validation, testing.

% ==================== SECTION 3 ====================
\section{All 17 Deep Learning Models}
\zhint{本课程涵盖的17个深度学习模型}

\begin{enumerate}[leftmargin=10pt,itemsep=0pt,topsep=0pt]
\item \textbf{MLP} -- basic classification/regression \zhint{多层感知机}
\item \textbf{DQN} -- basic RL agent \zhint{深度Q网络}
\item \textbf{DeepCoNN} -- recommender (review CNN) \zhint{深度协同神经网络}
\item \textbf{Wide \& Deep} -- recommender (memorize + generalize) \zhint{宽深网络}
\item \textbf{GCN} -- recommender (graph convolution) \zhint{图卷积网络}
\item \textbf{RNN} -- sequence modeling \zhint{循环神经网络}
\item \textbf{LSTM} -- sequence with gates \zhint{长短期记忆}
\item \textbf{BERT} -- encoder Transformer (NLP) \zhint{双向编码器}
\item \textbf{GPT} -- decoder Transformer (text gen) \zhint{生成预训练Transformer}
\item \textbf{Tacotron 2.0} -- speech synthesis (TTS) \zhint{语音合成}
\item \textbf{HuBERT} -- speech recognition (STT) \zhint{语音识别}
\item \textbf{CNN} -- image captioning encoder \zhint{卷积神经网络}
\item \textbf{GAN} -- image generation \zhint{生成对抗网络}
\item \textbf{Diffusion Model} -- image generation \zhint{扩散模型}
\item \textbf{D3QN} -- game agent \zhint{对决双重DQN}
\item \textbf{A3C} -- game agent (actor-critic) \zhint{异步优势Actor-Critic}
\item \textbf{AlphaZero} -- game agent (MCTS + self-play) \zhint{自我对弈}
\end{enumerate}

Plus: \textbf{TripoSR} (2D$\to$3D) and \textbf{Shap-E} (text$\to$3D) for 3D generation.

% ==================== SECTION 4 ====================
\section{Chapter-by-Chapter Quick Reference}
\zhint{各章快速索引}

\subsection{Ch2--5: Foundations}
\textbf{Data types}: structured (tabular) vs unstructured (text/image/audio).\\
\textbf{Preprocessing}: missing values (drop/impute), encoding (one-hot, label), scaling (StandardScaler, MinMaxScaler), PCA for dim reduction.\\
\textbf{Data representation}: TF-IDF (text), embeddings (word2vec), spectrograms (audio), pixel arrays (image).

\subsection{Ch6: ML -- Classification \& Regression}
\textbf{Classification}: DT, RF, SVM, KNN, ANN. Metrics: accuracy, precision, recall, F1, MCC, AUC-ROC.\\
\textbf{Regression}: Linear, Logistic, Ridge, LASSO, SVR, ANN. Metrics: $R^2$, MAE, RMSE.\\
\textbf{Model selection}: k-fold CV, train/val/test split.

\subsection{Ch7: Clustering \& RL}
\textbf{Clustering}: K-Means (elbow method), Agglomerative (linkage types), Spectral.\\
\textbf{Metrics}: Silhouette $[-1,1]\uparrow$, Calinski-Harabasz $\uparrow$, Davies-Bouldin $\downarrow$.\\
\textbf{RL}: agent/env/state/action/reward, Q-Learning, $\varepsilon$-greedy.

\subsection{Ch8: Data Visualization}
\textbf{Preprocessing}: histogram, boxplot, scatter, Q-Q, heatmap, pairplot.\\
\textbf{Analysis}: ROC, training curve, residual, calibration, learning curve.\\
\textbf{Combining}: subplots, overlay, shared axes. Text: word cloud.\\
\textbf{7 best practices}: audience, simplicity, color, labels, scale, context, honesty.

\subsection{Ch9: Business -- RecSys \& Stock}
\textbf{RecSys 3 generations}:\\
1st: CF (user/item), content-based, knowledge-based.\\
2nd: hybrid, matrix factorization ($R \approx PQ^T$).\\
3rd: DeepCoNN, Wide\&Deep, GCN.\\
\textbf{Stock}: yfinance, EMA, LSTM (sliding window), DQN trading.\\
\textbf{Metrics}: Sharpe ratio, max drawdown, Sortino ratio.

\subsection{Ch10: NLP -- Semantic \& Text Gen}
\textbf{Semantic}: TF-IDF+NB (0.84), NBoW (0.86), LSTM, BERT (MLM+NSP), Transformer (TinyBERT 0.86).\\
\textbf{BERT variants}: ALBERT, RoBERTa, DistilBERT, TinyBERT.\\
\textbf{Benchmarks}: GLUE (9 tasks), SQuAD (QA).\\
\textbf{Text gen}: Markov chain, GPT (1$\to$4), LLaMA.\\
\textbf{Architecture}: encoder-only (BERT) vs decoder-only (GPT) vs seq2seq (T5).

\subsection{Ch11: Multimedia}
\textbf{Audio}: vocal separation (librosa, Demucs).\\
\textbf{TTS}: SPSS, concatenative, WaveNet (dilated causal conv), Tacotron2. Metric: MOS.\\
\textbf{STT}: Wav2Vec 2.0, HuBERT (k-means pseudo-labels).\\
\textbf{Image caption}: CNN-RNN+attention, ViT+Transformer. Metrics: BLEU, METEOR, CIDEr, SPICE.\\
\textbf{Image gen}: GAN, DALL-E+CLIP, diffusion. Metric: FID.\\
\textbf{Video}: DeepFake (encoder-decoder swap).

\subsection{Ch12: Game}
\textbf{Blackjack}: state=(sum, dealer, ace), hit/stand.\\
\textbf{D3QN}: double (reduce overestimation) + dueling ($V(s)+A(s,a)$).\\
\textbf{A3C}: async parallel workers, actor-critic, advantage.\\
\textbf{AlphaZero}: MCTS + dual-headed net + self-play.\\
\textbf{EWC}: prevent catastrophic forgetting across games.\\
\textbf{3D gen}: TripoSR (image$\to$3D), Shap-E (text$\to$3D).\\
\textbf{Robot}: transfer learning, imitation, behavior cloning, inverse RL, diffusion policies.

% ==================== SECTION 5 ====================
\section{Future of Data Science}
\zhint{数据科学的未来}

\textbf{Roles}: Data Scientist (models) vs ML Engineer (deploy) vs Data Engineer (pipeline/infra).\\
\textbf{Emerging}: IoT (smart sensors), spatial intelligence (3D reasoning), Vision-Language Models (VLMs, multimodal).\\
\textbf{Final message}: embrace GenAI tools; think products over applications; specialized agents over general ones.

\keybox{
\textbf{Course in One Line}:\\
Data pipeline (collect $\to$ preprocess $\to$ model $\to$ interpret) applied to business, NLP, multimedia, and games using 17 deep learning models from MLP to diffusion.\\
\zhint{数据流水线应用于商业/NLP/多媒体/游戏，涵盖17个深度学习模型}
}

\end{multicols}
\end{CJK}
\end{document}
